\documentclass[../main.tex]{subfiles}

\begin{document}
\subsection{Sucesiones de funciones}
\classdate{2026-05-06}

\begin{definition}[Convergencia puntual]
    Sea \(\set{f_{n} \colon U \subseteq \mathbb{C} \to \mathbb{C} \given f\;
    \text{función}}_{n\in \mathbb{N}}\) una sucesión de funciones. Decimos que
    \(\{f_{n}\}\) converge putualmente a \(f\) si \(\forall \varepsilon > 0\) y
    \(\forall z\in U \exists N\in \mathbb{N}\) tal que \(\forall n \geq N\),

    \begin{equation*}
        \abs{f_{n}(z) - f(z)} < \varepsilon.
    \end{equation*}
\end{definition}

\begin{definition}[Convergencia uniforme]
    Sea \(\{f_{n} \colon U \subseteq \mathbb{C} \to \mathbb{C}\}_{n\in \mathbb{N}}\) una
    sucesión de funciones. Decimos que \(\{f_{n}\}\) converge uniformemente a \(f\) si
    \(\forall \varepsilon > 0\) existe \(N\in \mathbb{N}\) tal que \(n \geq N\)
    \(\abs{f_{n}(z) - f(z)} < \varepsilon\; \forall z\in U\).
\end{definition}

\begin{proposition}\label{prop:convergenciaLimit}
    Sea \(\{f_{n} \colon U \subseteq \mathbb{C} \to \mathbb{C}\}_{n\in \mathbb{N}}\) una
    sucesión de funciones tal que \(\{f_{n}\}\) converge uniformemente a \(f \colon U
    \subseteq \mathbb{C} \to \mathbb{C}\) y sea \(\gamma \subseteq U\) una curva suave
    por pedazos, entonces

    \begin{enumerate}
        \item Si \(\forall n\in \mathbb{N}\) \(f_{n}\) es continua, entonces
            \(f\) es continua.
        \item \(\lim_{n \to \infty} \int_{\gamma}^{}f_{n}(z)\odif{z} =
                \int_{\gamma}^{}\lim_{n \to \infty} f_{n}(z)\odif{z} =
            \int_{\gamma}^{}f(z)\odif{z}\) con \(f_{n}\) continuas.
    \end{enumerate}
\end{proposition}

\begin{proof}
    Debemos mostrar que \(\forall z_{0}\in U\) y \(\forall \varepsilon > 0 \;\exists
    \delta > 0\) tal que \(\abs{z - z_{0}} < \delta\), entonces \(\abs{f(z) - f(z_{0})}
    < \varepsilon\). Sea \(\varepsilon > 0\) y \(z_{0}\in U\), \(\exists N\in
    \mathbb{N}\) tal que \(\forall n \geq N \abs{f_{n}(z_{0}) - f(z_{0})} <
    \tfrac{\varepsilon}{3}\), en particular \(\abs{f_{N}(z) - f(z_{0})} <
    \tfrac{\varepsilon}{3}\).

    Como \(f_{N}\) es continua \(\exists\delta > 0\) tal que cumple
    continuidad. Sea \(z\in
    U\) tal que \(\abs{z - z_{0}}  \delta\) se cumple que \(\abs{f_{N}(z) -
        f_{N}(z_{0})} <
    \tfrac{\varepsilon}{3}\).
    Y nuevamente por convergencia uniforme,

    \begin{equation*}
        \abs{f_{N}(z) - f(z)} < \dfrac{\varepsilon}{3}.
    \end{equation*}

    Notemos que \(z\in U\) tal que \(\abs{z - z_{0}} < \delta\) y
    \begin{align*}
        \abs{f(z) - f(z_{0})} \leq \abs{f(z) - f_{N}(z)} + \abs{f_{N}(z) - f_{N}(z_{0})}
        + \abs{f_{N}(z_{0}) - f(z_{0})} \leq \dfrac{\varepsilon}{3} +
        \dfrac{\varepsilon}{3} + \dfrac{\varepsilon}{3} = \varepsilon.
    \end{align*}

    Por lo tanto, \(f\) es continua en \(z_{0},\; \forall z_{0}\in U\).

    \classdate{2026-05-07}

    Para la segunda, queremos mostrar que \(\lim_{n \to \infty}
    \int_{\gamma}^{}f_{n}(z)\odif{z} = \int_{\gamma}^{}f(z)\odif{z}\). Es decir, que
    \(\forall \varepsilon > 0\; \exists N\in \mathbb{N}\) tal que \( \forall n \geq N\),

    \begin{equation*}
        \abs`\Bigg{\int_{\gamma}^{}f_{n}(z)\odif{z} - \int_{\gamma}^{}f(z)\odif{z}}
        < \varepsilon.
    \end{equation*}

    Sea \(\varepsilon > 0\). Como \(f_{n}\) converge uniformemente a \(f\), i.e. \(f_{n}
    \xlongrightarrow[]{u} f (f_{n} \rightrightarrows f)\) existe \(N\in \mathbb{N}\;
    \forall n \geq N\)

    \begin{equation*}
        \abs{f_{n}(z) - f(z)} <
        \dfrac{\varepsilon}{\mathrm{long}(\gamma)}\quad \forall z\in
        U.
    \end{equation*}

    Entonces

    \begin{align*}
        \abs`\Bigg{\int_{\gamma}^{}f_{n}(z)\odif{z} - \int_{\gamma}^{}f(z)\odif{z}} &=
        \abs`\Bigg{\int_{\gamma}^{}f_{n}(z) - f(z)\odif{z}},\\
        &\leq \int_{\gamma}^{}\abs{f_{n}(z) - f(z)}\abs{\odif{z}},\\
        &\leq
        \dfrac{\varepsilon}{\mathrm{long}(\gamma)}\int_{\gamma}^{}\abs{\odif{z}},\\
        &= \dfrac{\varepsilon}{\mathrm{long}(\gamma)}\mathrm{long}(\gamma),\\
        &= \varepsilon.
    \end{align*}
\end{proof}

\begin{proposition}
    Sea \(\{f_{n} \colon U \subseteq \mathbb{C} \to \mathbb{C}\}_{n\in \mathbb{N}}\) una
    sucesión de funciones que converge uniformemente a \(f \colon U \subseteq \mathbb{C}
    \to \mathbb{C}\) sobre cualquier bola cerrada contenida en \(U\). Si \(f_{n}\) es
    analítica \(\forall n\in \mathbb{N}\), entonces \(f\) es analítica y mas aún para
    cualquier bola cerrada

    \begin{equation*}
        \lim_{n \to \infty} f_{n}^{(k)}(z) = f^{(k)}(z)\qquad \forall z\in U.
    \end{equation*}
\end{proposition}

\begin{proof}
    Como \(\forall n\in \mathbb{N}\) \(f_{n}\) es analítica, entonces \(\forall n\in
    \mathbb{N}\) \(f_{n}\) es continua, entonces por (a) de
    \zcref{prop:convergenciaLimit}, \(f\) es
    continua y sea \(\gamma \subseteq U\) una curva cerrada suave por pedazos. Como
    \(f_{n}\) es analítica \(\forall n\in \mathbb{N}\),

    \begin{equation}
        \int_{\gamma}^{}f_{n}(z)\odif{z} = 0.
        \label{eq:curvaCerradaCero}
    \end{equation}

    Ahora, por (b) de la \zcref{prop:convergenciaLimit},

    \begin{equation*}
        \lim_{n \to \infty} \int_{\gamma}^{}f_{n}(z)\odif{z} =
        \int_{\gamma}^{}f(z)\odif{z}
    \end{equation*}

    pero por \zcref[noname]{eq:curvaCerradaCero},

    \begin{equation*}
        \int_{\gamma}^{}f(z)\odif{z} = 0
    \end{equation*}

    \(\forall \gamma\in U \) curva cerrada suave por pedazos. Entonces por el
    \zcref{theorem:Morera}, \(f\) es analítica en
    \(U\).

    Sea \(z_{0}\in  U\), \(r > 0\) ``pequeña'' tal que \(\bar{B}_{r}(z_{0}) \subseteq
    U\) y \(R > 0\) tal que \(\partial B_{R}(z_{0}) \subseteq U\).
    Sea \(\gamma \colon [0, 2\pi] \to U\) una parametrización en sentido antihorario de
    \(\partial B_{R}(z_{0})\) que la recorra solo una vez.

    Sea \( z\in \bar{B}_{r}(z_{0})\), entonces

    \begin{align*}
        f_{n}^{(k)}(z) &= \dfrac{k!}{2\pi i}\int_{\gamma}^{}\dfrac{f_{n}(w)}{\bigl(w -
        z\bigr)^{k + 1}}\odif{w},\\
        f^{(k)}(z) &= \dfrac{k!}{2\pi i}\int_{\gamma}^{}\dfrac{f(w)}{\bigl(w
            - z\bigr)^{k +
        1}}\odif{w}.
    \end{align*}

    Sea \(\varepsilon > 0\), entonces existe \(N\in \mathbb{N}\) tal que
    \(\forall n \geq
    N\),

    \begin{equation*}
        \abs{f_{n}(z) - f(z)} < \varepsilon\quad \forall z\in \bar{B}_{R}(z_{0}).
    \end{equation*}

    Notemos que

    \begin{align*}
        \abs{f_{n}^{(k)}(z) - f^{(k)}(z)} &=
        \dfrac{k!}{2\pi}\abs`\Bigg{\int_{\gamma}^{}\dfrac{f_{n}(w) -
            f(w)}{\bigl(w - z\bigr)^{k +
        1}}\odif{w}},\\
        &\leq \dfrac{k!}{2\pi}\int_{\gamma}^{}\dfrac{\abs{f_{n}(w) - f(w)}}{\abs{w -
        z}^{k + 1}}\abs{\odif{w}},\\
        &< \dfrac{k!}{2\pi}\int_{\gamma}^{}\dfrac{\varepsilon}{\abs{w - z}^{k +
        1}}\abs{\odif{w}},\\
        \intertext{Además que \(\forall w\in \partial B_{R}(z_{0}), \forall z\in
            \bar{B}_{r}(z_{0})\) se tiene que \(\abs{w - z} \geq R - r\), lo
            cual implica
        que \(\tfrac{1}{\abs{w - z}} \leq \tfrac{1}{R - r}\), entonces}
        &\leq \dfrac{k!}{2\pi}\int_{\gamma}^{}\dfrac{\varepsilon}{\bigl(R - r\bigr)^{k +
        1}}\abs{\odif{w}},\\
        &= \dfrac{k!}{2\pi}\dfrac{\varepsilon}{\bigl(R - r\bigr)^{k +
        1}}\int_{\gamma}^{}\abs{\odif{w}},\\
        &= \dfrac{k! R}{\bigl(R - r\bigr)^{k + 1}}\varepsilon.
    \end{align*}

    Por lo que,

    \begin{equation*}
        \lim_{n \to \infty} f_{n}^{(k)}(z) = f^{(k)}(z),\medspace \forall z\in
        \bar{B}_{r}(z_{0}).
    \end{equation*}
\end{proof}

\begin{definition}
    Decimos que una serie de funciones \(\sum\limits_{n = 1}^{\infty}f_{n}\) converge
    uniformemente a \(f\), si la sucesión de funciones \(g_{n} = \sum\limits_{i =
    1}^{n}f_{i}\) converge uniformemente a \(f\).
\end{definition}

\classdate{2026-05-12}

\begin{theorem}[Criterio de Cauchy]
    Sean \(A \subseteq \mathbb{C}\), \(f_{n} \colon A \subseteq \mathbb{C} \to
    \mathbb{C}\) una sucesión de funciones. Entonces \(\exists f \colon A \subseteq
    \mathbb{C} \to \mathbb{C}\) tal que \(f_{n}
    \xlongrightarrow[]{u} f\;\) en \(A\) si y solo si \(\forall \varepsilon > 0,\;
    \exists N\in \mathbb{N}\) tal que si \(m, n \geq N\), entonces \(\abs{f_{n}(z) -
    f_{m}(z)} < \varepsilon\; \forall z\in A\). (\(f_{n}\) es uniformemente de Cauchy).
\end{theorem}

\begin{proof}
    La primera implicación es directa. Por lo que probaremos que si la sucesión
    \(f_{n}\) es uniformemente de Cauchy implica que converge uniformemente a \(f\) A.

    Sea \(\varepsilon > 0\). Por hipótesis \(\exists N\in \mathbb{N}\) tal que si \(n,m
    \geq N\), entonces \(\abs{f_{n}(z) - f_{m}(z)} <
        \tfrac{\varepsilon}{2},\; \forall z\in
    A\).
    En particular, \(\forall z\in A\), \(\{f_{n}(z)\}\) es de Cauchy. Como
    \(\mathbb{C}\) es completo, \(\exists f(z)\in \mathbb{C}\) tal que \(\lim_{n \to
    \infty} f_{n}(z) = f(z)\). Por lo que definimos \(f \colon A \subseteq \mathbb{C}
    \to \mathbb{C}\) como \(f(z) \coloneqq \lim_{n \to \infty} f_{n}(z)\).

    \textcolor{JungleGreen}{Afirmación.} \(f_{n}
    \xlongrightarrow[]{u} f\) en \(A\).

    Como \(f(z) \coloneqq \lim_{n \to \infty} f_{n}(z)\), para cada \(z\in A,\; \exists
    m_{z}\in \mathbb{N}\) tal que \(\abs{f(z) - f_{n + m_{z}}(z)} <
    \tfrac{\varepsilon}{2}\).
    Sea \(z\in A,\; n \geq N\),

    \begin{align*}
        \abs{f(z) - f_{n}(z)} &= \abs{f(z) - f_{N + m_{z}}(z) + f_{N + m_{z}}(z) -
        f_{n}(z)},\\
        &\leq \abs{f(z) - f_{N + m_{z}}(z)} + \abs{f_{N + m_{z}}(z) -
        f_{n}(z)},\\
        &< \dfrac{\varepsilon}{2} + \dfrac{\varepsilon}{2} =
        \varepsilon.
    \end{align*}

    Por lo tanto, \(f_{n} \xlongrightarrow[]{u} f \) en \(A\).
\end{proof}

\begin{remark}
    Si \(f_{k} \colon A \subseteq \mathbb{C} \to \mathbb{C}\) es tal que \(f_{n}
    \xlongrightarrow[]{u} f\), con \(f \colon A \subseteq \mathbb{C} \to \mathbb{C}\) y
    \(f_{k}\) es continua \(\forall k\in \mathbb{N}\), entonces \(f\) es continua.
\end{remark}

\begin{corollary}
    Si \(f_{k} \colon A \subseteq \mathbb{C} \to \mathbb{C}\) es una sucesión de
    funciones continuas tales que \(\sum_{k = 1}^{\infty}f_{k}\) converge uniformemente,
    entonces \(\sum_{k = 1}^{\infty}f_{k}\) es continua.
\end{corollary}

\begin{definition}
    Sea \(f_{k} \colon A \subseteq \mathbb{C} \to \mathbb{C}\) una sucesión de
    funciones. Decimos que \(\sum_{k = 1}^{\infty}f_{k}\) converge absoluta y
    uniformemente en \(A\), si \(\sum_{k = 1}^{\infty}\abs{f_{k}}\) converge
    uniformemente.
\end{definition}

\begin{remark}
    Por el criterio uniformemente de Cauchy,

    \begin{equation*}
        \abs`\Bigg{\sum_{k = n + 1}^{n + m}f_{k}(z)} \leq \sum_{k = n + 1}^{n +
        m}\abs{f_{k}(z)},
    \end{equation*}

    \(\forall z\in A,\; \forall n,m\in \mathbb{N}\).
    Si \(\sum_{k = 1}^{\infty}f_{k}(z)\) converge absoluta y uniformemente, entonces
    \(\sum_{k = 1}^{\infty}f_{k}\) converge uniformemente.
\end{remark}

\begin{theorem}[Prueba M de Weierstrass]\label{thm:MTestWeierstrass}
    Sean \(A \subseteq \mathbb{C}\) y \(f_{n} \colon A \subseteq \mathbb{C} \to
    \mathbb{C}\) una sucesión de funciones. Supongamos que existen \(M_{n} \geq 0\),
    \(n\in \mathbb{N}\) tales que

    \begin{enumerate}
        \item \(\abs{f_{n}(z)} \leq M_{n},\; \forall z\in A\).
        \item \(\sum_{n = 1}^{\infty}M_{n}\) converge, \(\sum_{n = 1}^{\infty}f_{n}\)
            converge absoluta y uniformemente en \(A\).
    \end{enumerate}
\end{theorem}

\begin{proof}
    Como \(\sum_{n = 1}^{\infty}M_{n}\) converge, dada \(\varepsilon > 0, \;
        \exists N\in
    \mathbb{N}\) tal que \(\forall m\in \mathbb{N}\), \(\forall m \geq N\),

    \begin{equation*}
        \abs`\Bigg{\sum_{k = n + 2}^{n + m}M_{k}} < \varepsilon.
    \end{equation*}

    Sean \(n \geq N, \; m\in \mathbb{N}\),

    \begin{align*}
        \sum_{k = n + 1}^{n + m}\abs{f_{k}(z)} &\leq \sum_{k = n + 1}^{n + m}M_{k},\\
        &= \abs`\Bigg{\sum_{k = n + 1}^{n + m}M_{k}} < \varepsilon.
    \end{align*}

    Por lo que \(\sum_{k = 1}^{\infty}\abs{f_{k}}\) es uniformemente de Cauchy y, por lo
    tanto, \(\sum_{k = 1}^{\infty}f_{k}\) converge absoluta y uniformemente en \(A\).
\end{proof}

\begin{example}
    Consideremos \(\sum_{n = 1}^{\infty}\frac{z^{n}}{n}\).

    Notemos que esta serie converge absoluta y uniformemente en
    \(\overline{B_{r}(0)}\) con
    \(r < 1\).
    Sea \(f_{n}(z) = \frac{z^{n}}{n}\),

    \begin{equation*}
        \abs{f_{n}(z)} = \dfrac{\abs{z}^{n}}{n} \leq \dfrac{r^{n}}{n} \leq r^{n}.
    \end{equation*}

    Como \(0 < r < 1\), \(\sum_{n = 1}^{\infty}r^{n}\) converge. Por el
    \zcref{thm:MTestWeierstrass}, \(\sum_{n =
    1}^{\infty}f_{n}\) converge absoluta y uniformemente en \(\overline{B_{r}(0)}\).
\end{example}

\classdate{2026-05-14}

\begin{theorem}[de Weierstrass]\label{thm:deWeierstrass}
    Sea \(\{f_{n} \colon U \subseteq \mathbb{C} \to \mathbb{C}\}\) una sucesión de
    funciones.

    \begin{enumerate}
        \item Si \(f_{n} \xlongrightarrow[]{u} f\) en bolas cerrada y
            \(f_{n}\) es analítica
            \(\forall n\in \mathbb{N}\), entonces \(f\) es analítica y
            \(\lim_{n \to \infty}
            f_{n}^{(k)}(z) = f^{(k)}(z)\; \forall k\in \mathbb{N}\) y \(\forall z\in
            \overline{B}\).
        \item Si \(\sum_{n = 0}^{\infty}f_{n}\) converge uniformemente a \(f\) en bolas
            cerraddas y \(f_{n}\) es analítica \(\forall n\in \mathbb{N}\),
            entonces \(f\)
            es analítica y \(\sum_{n = 0}^{\infty}f_{n}^{(k)}(z) =
                f^{(k)}(z)\; \forall k\in
            \mathbb{N}\) y \(\forall z\in \overline{B}\).
    \end{enumerate}
\end{theorem}

\begin{definition}
    Decimos que \(\{f_{n} \colon U \subseteq \mathbb{C} \to \mathbb{C}\}\) converge
    normalmente a \(f \colon U \subseteq \mathbb{C} \to \mathbb{C}\) si \(f_{n}\)
    converge uniformemente a \(f\; \forall k \subseteq U\) compacto. Y si \( U\) es
    compacto, la convergencia normal implica convergencia uniforme.
\end{definition}

\begin{example}
    Sea \(f(z) = \sum_{n = 0}^{\infty}\tfrac{z^{n}}{e^{n}}\) con \(f \colon
        \bar{B}_{1}(0)
    \to \mathbb{C}\).

    Notemos que

    \begin{equation*}
        \abs{\dfrac{z^{n}}{e^{n}}} \leq \dfrac{1}{e^{n}}
    \end{equation*}

    Recordando que \(n^{2} \leq e^{n}\; \forall n\in \mathbb{N}\), tal que
    \(\dfrac{1}{e^{n}} \leq \dfrac{1}{n^{2}}\) y como \(\sum_{n =
    1}^{\infty}\tfrac{1}{n^{2}}\) converge, \(\sum_{n = 0}^{\infty}\tfrac{1}{e^{n}}\)
    también converge.
    Y por la prueba M de Weierstrass, \(\sum_{n =
    0}^{\infty}\tfrac{z^{n}}{e^{n}}\) converge
    uniforme y absolutamente, además por el \zcref{thm:deWeierstrass}
    \(f\) es analítica en \(\bar{B}_{1}(0)\) y, además,

    \begin{equation*}
        f^{\prime}(z) = \sum_{n = 0}^{\infty}\dfrac{n z^{n - 1}}{e^{n}}.
    \end{equation*}
\end{example}

\begin{definition}
    Sea \(\{f_{n} \colon \mathbb{C} \to \mathbb{C}\}\) una sucesión de funciones tal que
    \(\exists z_{0}\in \mathbb{C}\) y \(\forall n\in \mathbb{N}\; \exists a_{n}\in
    \mathbb{C}\) tal que \(f_{n}(z) = a_{n}\bigl(z - z_{0}\bigr)^{n}\). En cuyo caso a
    la sucesión \(\{f_{n}\}\) se le llama sucesión de potencias y a su respectiva serie

    \begin{equation*}
        f(z) = \sum_{n = 0}^{\infty}f_{n}(z) = \sum_{n = 0}^{\infty}a_{n}(z - z_{0})^{n}
    \end{equation*}

    se le llama serie de potencias.
\end{definition}

\begin{remark}
    Recordemos ahora que significaba que una función fuese analítica en
    \(\mathbb{R}\). Es
    decir, sea \(f \colon U \subseteq \mathbb{R} \to \mathbb{R}\) analítica si \(\forall
    x_{0}\in U\; \exists \delta > 0\) tal que

    \begin{equation*}
        f(x) = \sum_{n = 0}^{\infty}\dfrac{f^{(n)}(x_{0})}{n!}(x -
        x_{0})^{n}\quad \forall
        x\in (x_{0} - \delta, x_{0} + \delta).
    \end{equation*}

    En \(\mathbb{C}\) sabemos que si una función es derivable y satisface las
    ecuaciones de
    Cauchy-Riemman la función analítica. Sin embargo, para el caso real, esto
    no siempre el
    caso. Por lo que veremos una función que es de clase \(C^{\infty}\), pero no es
    analítica.

    \begin{example}
        Sea

        \begin{equation*}
            f(x) =
            \begin{dcases}
                \mathrm{e}^{-\tfrac{1}{x^{2}}} & x \neq 0,\\
                0 & x = 0.
            \end{dcases}
        \end{equation*}

        pues \(f^{(k)}(0) = 0,\; \forall k\in \mathbb{N}\), lo cual implica que

        \begin{equation}
            g(x) = \sum_{n = 0}^{\infty}\dfrac{f^{(k)}(0)}{n!}x^{n} = 0,\;
            \forall \delta > 0,
        \end{equation}

        por lo que \(f(x) \neq g(x) = 0\; \forall x\in (-\delta, \delta)\).
    \end{example}
\end{remark}

\begin{lemma}[Abel]\label{lemma:abel}
    Si existe \(z_{1}\in \mathbb{C}\) tal que la serie \(\sum_{n =
    0}^{\infty}a_{n}(z_{1} - z_{0}) ^{n}\) converge, entonces la serie \(\sum_{n =
    0}^{\infty}a_{n}(z - z_{0})^{n}\) converge normal y absolutamente \(\forall z\in
    B_{R}(z_{0})\), donde \(R = \abs{z_{1} - z_{0}}\).
\end{lemma}

\begin{proof}
    Sea \(R = \abs{z_{1} - z_{0}}\). Es claro que \(\abs{a_{n}}R^{n} \xlongrightarrow[n
    \to \infty]{} 0\) y, por tanto, está acotada por algún \(M > 0\). Ahora, sea \(r <
    R\) y sea la sucesión

    \begin{equation*}
        M_{n} = M\bigl(\dfrac{r}{R}\bigr)^{n}.
    \end{equation*}

    Como \(\tfrac{r}{R} < 1\), la serie geométrica \(\sum_{n =
    0}^{\infty}M(\tfrac{r}{R})^{n}\) converge, i.e. \(\sum_{n =
    0}^{\infty}M_{n}\) converge.
    Además en \(z\in \bar{B}_{r}(z_{0})\):

    \begin{align*}
        \abs{a_{n}(z_{1} - z_{0})^{n}} &\leq \abs{a_{n}}r^{n},\\
        &= \abs{a_{n}}\dfrac{R^{n}r^{n}}{R^{n}},\\
        &\leq M \dfrac{r^{n}}{R^{n}}.
    \end{align*}

    Y por la prueba M de Weierstrass concluimos que \(\sum_{n = 0}^{\infty}a_{n}(z -
    z_{0})^{n}\) converge uniforme y absolutamente para toda \(z\) en
    \(\bar{B}_{r}(z_{0})\) y \(r < R\). Por tanto, la serie converge normal y
    absolutamente
    para toda \(z\) en \(B_{r}(z_{0})\).
\end{proof}

\begin{proposition}[Radio de convergencia]\label{prop:radius-of-convergence}
    Sea \(\sum_{n = 0}^{\infty}a_{n}(z - z_{0})^{n}\) una serie de potencias, entonces
    \(\exists R\in [0, \infty]\)\footnote{El radio de convergencia es infinito, lo cual
    implica que la serie de potencias converge en todo \(\mathbb{C}\)} tal que si \(
    \abs{z - z_{0}} > R\) la serie diverge. Además, la serie es normal y absoluta en
    \(B_{R}(z_{0})\). A \(R\) se le llama el radio de convergencia de la
    serie de potencias.
\end{proposition}

\begin{proof}
    Sea \(R = \sup\set{r \geq 0 \given \sum_{n = 0}^{\infty}\abs{a_{n}}r^{n}\;
    \text{convere}}\). Sea \(r < R\), entonces la también converge, por lo que \(\sum_{n
    = 0}^{\infty}a_{n}(r_{1} + z_{0} - z_{0}) ^{n}\) converge.
    Usando el \zcref{lemma:abel} para \(r_{1} +
    z_{0} = z_{1}\), la serie \(\sum_{n = 0}^{\infty}a_{n}(z - z_{0})^{n}\) converge
    uniformemente para toda \(z\) en \(\bar{B}_{r_{1}}(z_{0})\) donde \(r_{1} =
    \abs{z_{1} - z_{0}} = \abs{r_{1} + z_{0} - z_{0}}\) y esto ocurre para toda \(r <
    r_{1} < R\). Por lo tanto, \(\forall z\in B_{R}(z_{0})\) converge normal y
    absolutamente.

    Ahora, suponemos que para \(z_{2}\in \mathbb{C}\) \(\abs{z_{2} - z_{0}} >
    R\), la serie
    \(\sum_{n = 0}^{\infty}a_{n}(z_{2} - z_{0})^{n} = \sum_{n =
    0}^{\infty}a_{n}s^{n}\), con
    \(s = \abs{z_{2} - z_{0}}\). De esta manera, notemos que \(\sum_{n =
        0}^{\infty}a_{n}(s
    + z_{0} - z_{0})^{n}\), con \(z_{2} = s + z_{0}\), la serie anterior
    converge para toda
    \(z\) en \(B_{s}(z_{0})\), lo cual es una contradicción pues \(s > R\).
\end{proof}

\classdate{2026-05-18}

\begin{theorem}[de los coeficientes de Taylor]
    Sea \(\sum_{n = 0}^{\infty}a_{n}(z - z_{0})^{n}\) una serie de potencias con un
    radio de convergencia \(R\) y definimos \(f \colon B_{R}(z_{0}) \subseteq \mathbb{C}
    \to \mathbb{C}\) dada por

    \begin{equation*}
        f(z) = \sum_{n = 0}^{\infty}a_{n}(z - z_{0})^{n}.
    \end{equation*}

    Entonces \(\forall n\in \mathbb{N}\)

    \begin{equation*}
        a_{n} = \dfrac{f^{(n)}(z_{0})}{n!}
    \end{equation*}

    y

    \begin{equation*}
        f^{\prime}(z_{0}) = \sum_{n = 0}^{\infty}n a_{n}(z - z_{0})^{n - 1}
    \end{equation*}

    tiene también a \(R\) como radio de convergencia.
\end{theorem}

\begin{remark}
    Una serie de potencias convergente \(f(z) = \sum a_{n}(z - z_{0})^{n}\) es analítica
    en el dominio de su radio de convergencia.
\end{remark}

\begin{proof}
    El \zcref{thm:deWeierstrass}afierma que
    \(f^{\prime}(z) = \sum_{n = 1}^{\infty}n a_{n}(z - z_{0})^{n - 1}\) para toda \(z\in
    \bar{B}_{r}(z_{0})\) y para toda \(z\in B_{R}(z_{0})\) con \(0 \leq r < R\).
    Suponemos que \(\exists z_{1}\in \mathbb{C}\) tal que \(\abs{z_{1} - z_{0}} > R\),
    por lo que \(\sum_{n = 1}^{\infty}n a_{n}(z - z_{0})^{n - 1}\) converge; tomando
    \(r_{1} = \abs{z_{1} - z_{0}}\), entonces \(\sum_{n = 1}^{\infty}n a_{n}r_{1}^{n -
    1}\), lo cual implica que \(n \abs{a_{n}}r_{1}^{n - 1} \to 0\). Por lo tanto,
    \(\exists M > 0\) \(n \abs{a_{n}}r_{1}^{n - 1} < M\) para toda \(n\) en
    los naturales, tal que

    \begin{equation*}
        \abs{a_{n}}r_{1}^{n} < \dfrac{M}{n} \leq M\; \forall n\in \mathbb{N}
    \end{equation*}

    y \(\abs{a_{n}}r_{1}^{n} \to 0\).

    Y, por el \zcref{lemma:abel}, \(\sum_{n =
    1}^{\infty}a_{n}r_{1}^{n}\) converge en \(\bar{B}_{r_{2}}(z_{0})\) con \(r_{1} >
    r_{2} > R\); sin embargo, esto nos lleva a una contradicción pues \(R\) es el
    radio de convergencia.

    Por lo tanto, el radio de convergencia de \(f^{\prime}\) debe de ser también \(R\).
\end{proof}

\begin{theorem}[Criterio de radio de convergencia]
    Sea \(\sum_{n = 0}^{\infty}a_{n}(z - z_{0})^{n}\) una serie de potencias y \(R\) el
    radio de convergencia.

    \begin{enumerate}
        \item Si \(\lim_{n \to \infty} \tfrac{\abs{n}}{\abs{a_{n + 1}}}\)
            existe y es \(L\),
            entonces el radio de convergencia es igual a este, i.e. \(R = L\).
        \item Si \(\lim_{n \to \infty} \sqrt[n]{\abs{a_{n}}}\) existe y es
            \(L\), entonces

            \begin{equation*}
                R =
                \begin{dcases}
                    \dfrac{1}{L} & \text{si}\; L \neq 0\; \wedge\; L \neq \infty,\\
                    0 & \text{si}\;  L = \infty,\\
                    \infty & \text{si}\; L = 0.
                \end{dcases}
            \end{equation*}
    \end{enumerate}
\end{theorem}

\begin{proof}
    Recordando que

    \begin{equation*}
        R = \sup\set{r \geq 0 \given \sum_{n = 0}^{\infty}\abs{a_{n}}r^{n}\;
        \text{convere}}
    \end{equation*}

    \begin{enumerate}
        \item Sea \(L = \lim_{n \to \infty} \tfrac{\abs{a_{n}}}{\abs{a_{n + 1}}}\) y
            suponiendo que \(0 < L < \infty\).
            Sea \(0 < r < L\), por lo que

            \begin{equation*}
                \lim_{n \to \infty} \dfrac{\abs{a_{n +
                1}}}{\abs{a_{n}}}\dfrac{r^{n + 1}}{r^{n}} =
                \dfrac{r}{L} < 1.
            \end{equation*}

            Por lo tanto, \(\sum_{n = 0}^{\infty}a_{n}r^{n}\) converge, en cambio si \(r
            > L\), lo cual implica que \(\tfrac{r}{L} > 1\), de tal forma que la serie
            anterior diverge, i.e. \(L = R\).

        \item Sea \(L = \lim_{n \to \infty} \sqrt[n]{\abs{a_{n}}}\). Suponiendo que
            \(0 < L < \infty\) y \(r > 0\). Notemos entonces que \(\lim_{n \to \infty}
                \sqrt[n]{\abs{a_{n}}r^{n}} = \lim_{n \to \infty}
            \sqrt[n]{\abs{a_{n}}}r = Lr\).
            Por el criterio de convergencia, si \(r > \tfrac{1}{L}\) la serie diverge y
            si \(r < \tfrac{1}{L}\), la serie converge.
            Por lo tanto, \(R = \tfrac{1}{L}\), tal que \(L = 0\) y \(L = \infty\).
    \end{enumerate}
\end{proof}

\begin{example}
    Sea \(\mathrm{e}^{z} = \sum_{n = 0}^{\infty} \tfrac{z^{n}}{n!}\), con \(a_{n} =
    \tfrac{1}{n!}\).

    \begin{equation*}
        \lim_{n \to \infty} \dfrac{\abs{a_{n}}}{\abs{a_{n + 1}}} = \lim_{n \to \infty}
        \dfrac{\tfrac{1}{n!}}{\tfrac{1}{(n + 1)!}} = \lim_{n \to \infty}
        \dfrac{(n + 1)!}{n!}
        = \lim_{n \to \infty} n + 1 = \infty.
    \end{equation*}

    Sea \(\sum_{n = 0}^{\infty}n^{n}z^{n}\) con \(a_{n} = n^{n}\), tal que

    \begin{equation*}
        \lim_{n \to \infty} \sqrt[n]{\abs{a_{n}}} = \lim_{n \to \infty} n = \infty
    \end{equation*}

    con \(R = 0\).
\end{example}

\pagebreak
\classdate{2026-05-19}
\begin{lemma}[1]\label{lemma:lem1}
    La serie de potencias \(\sum_{n = 0}^{\infty}z^{n}\) converge normal y absolutamente
    en \(B_{1}(0)\) a \(\tfrac{1}{1 - z}\).
\end{lemma}

\begin{proof}
    Recordamos que

    \begin{equation*}
        s_{n} = 1 + z + z^{2} + \cdots + z^{n} = \dfrac{1 - z^{n + 1}}{1 - z}.
    \end{equation*}

    Entonces

    \begin{align*}
        \sum_{n = 0}^{\infty}z^{n} &= \lim_{n \to \infty} s_{n} = \lim_{n \to
        \infty} 1 + z +
        \cdots + z^{n},\\
        &= \lim_{n \to \infty} \dfrac{1 - z^{n + 1}}{1 - z},
    \end{align*}

    como \(z\in B_{1}(0)\), i.e. \(\abs{z} < 1\) entonces

    \begin{equation*}
        \lim_{n \to \infty} \abs{z^{n + 1}} = \lim_{n \to \infty} \abs{z}^{n + 1} = 0.
    \end{equation*}

    Por lo que,

    \begin{equation*}
        \lim_{n \to \infty} \dfrac{1 - z^{n + 1}}{1 - z} = \lim_{n \to
        \infty} \dfrac{1}{1 -
        z} - \lim_{n \to \infty} \dfrac{z^{n + 1}}{1 - z} = \dfrac{1}{1 - z}.
    \end{equation*}

    Por lo tanto,

    \begin{equation*}
        \sum_{n = 0}^{\infty}z^{n} = \dfrac{1}{1 - z}.
    \end{equation*}
\end{proof}

\begin{lemma}[2]\label{lemma:lem2}
    Si \(\sum_{n = 0}^{\infty}f_{n}\) converge uniformemente a \(f\) (\(f_{n},\; f
    \colon U \subseteq \mathbb{C} \to \mathbb{C}\; \forall n\in \mathbb{N}\)) y \(g
    \colon U \subseteq \mathbb{C} \to \mathbb{C}\) acotada. Entonces \(\sum_{n =
    0}^{\infty}g f_{n}\) converge uniformemente a \(gf\).
\end{lemma}

\begin{proof}
    Como \(g\) está acotada \(\exists n > 0\) tal que \(\abs{g(z)} < M\) para toda \(z\)
    en \(U\). Y, además, la serie \(\sum_{n = 0}^{\infty}f_{n}\) converge uniformemente
    y es uniformemente de Cauchy.

    Sea \(\varepsilon > 0\) entonces \(\exists N\in \mathbb{N}\) tal que
    \(\forall m \geq k
    > N\),

    \begin{equation*}
        \abs`\bigg{\sum_{n = k}^{m}f_{n}(z)} < \dfrac{\varepsilon}{M}\; \forall z\in U.
    \end{equation*}

    Notemos que

    \begin{align*}
        \abs`\bigg{\sum_{n = k}^{m}g(z)f_{n}(z)} &\leq \abs`\bigg{g(z)\sum_{n
        = k}^{m}f_{n}(z)},\\
        &= \abs{g(z)}\abs`\bigg{\sum_{n = k}^{m}f_{n}(z)},\\
        &\leq M \abs`\bigg{\sum_{n = k}^{m}f_{n}(z)},\\
        &< M\dfrac{\varepsilon}{M} = \varepsilon.
    \end{align*}
\end{proof}

\begin{theorem}[de Taylor]
    Sea \(f \colon U \subseteq \mathbb{C} \to \mathbb{C}\) analítica, entonces \(\exists
    r > 0\) tal que \(\bar{B}_{r}(z_{0}) \subseteq U\) y \(\forall z\in B_{r}(z_{0})\),

    \begin{equation*}
        f(z) = \sum_{n = 0}^{\infty}\dfrac{f^{(n)}(z_{0})}{n!}(z - z_{0})^{n}.
    \end{equation*}
\end{theorem}

\begin{proof}
    Como \(U\) es abierto dado \(z_{0}\in  U\) \(\exists r > 0\) ``pequeño'' tal que
    \(\bar{B}_{r}(z_{0}) \subseteq U\). Sea \(t > r\) tal que \(\bar{B}_{t}(z_{0})\) y
    \(\gamma \colon [0, 2\pi] \to \mathbb{C}\) dada por \(\gamma(s) = z_{0} + t
    \mathrm{e}^{is}\), \(\gamma = C_{t}(z_{0})\). Por la fórmula integral de Cauchy, si
    \(z\in \bar{B}_{r}(z_{0})\) entonces

    \begin{equation*}
        f(z) = \dfrac{1}{2\pi i} \int_{\gamma}^{}\dfrac{f(z)}{w - z}\odif{w}
    \end{equation*}

    con \(\abs{z - z_{0}} \leq r < t\) y, además, \(w\in \gamma\) tal que \(\abs{w -
    z_{0}}\),

    \begin{equation*}
        \abs{\dfrac{z - z_{0}}{w - z_{0}}} = \dfrac{\abs{z - z_{0}}}{\abs{w -
        z_{0}}} \leq
        \dfrac{r}{t} < 1.
    \end{equation*}

    Notemos que

    \begin{align*}
        \dfrac{1}{w - z} &= \dfrac{1}{w - z_{0} - (z - z_{0})},\\
        &= \dfrac{1}{w - z_{0}}\dfrac{1}{1 - \tfrac{z - z_{0}}{w - z_{0}}},\\
        \intertext{Por el \zcref{lemma:lem1},}
        &= \dfrac{1}{w - z_{0}}\sum_{n = 0}^{\infty}\bigl(\tfrac{z - z_{0}}{w -
        z_{0}}\bigr)^{n}.
    \end{align*}

    Ahora, sea \(g \colon \gamma \to U\) dada por \(g(w) = \tfrac{f(w)}{w - z_{0}}\) con
    \(g\) continua y \(\gamma\) compacta, lo cual implica que \(g\) es compacta.
    Y dado que
    \begin{equation*}
        \sum_{n = 0}^{\infty}\bigl(\tfrac{z - z_{0}}{w - z_{0}}\bigr)^{n}
    \end{equation*}

    converge uniformemente, entonces

    \begin{equation*}
        \sum_{n = 0}^{\infty}\dfrac{f(w)(z - z_{0})^{n}}{(w - z_{0})^{n + 1}}
    \end{equation*}

    converge. Por lo tanto,

    \begin{equation*}
        \dfrac{f(w)}{w - z_{0}}\dfrac{1}{1 - \tfrac{z - z_{0}}{w - z_{0}}} =
        \dfrac{f(w)}{w -
        z}.
    \end{equation*}

    Por convergencia uniforme,

    \begin{equation*}
        \int_{\gamma}^{}\dfrac{f(w)}{w - z}\odif{w} = \int_{\gamma}^{}\sum_{n =
        0}^{\infty}\dfrac{f(z)(z - z_{0})^{n}}{(w - z_{0})^{n + 1}}\odif{w} = \sum_{n =
        0}^{\infty}\int_{\gamma}^{}\dfrac{f(z)(z - z_{0})^{n}}{(w - z_{0})^{n
        + 1}}\odif{w}.
    \end{equation*}

    Por lo que

    \begin{align*}
        f(z) &= \dfrac{1}{2\pi i}\int_{\gamma}^{}\dfrac{f(z)}{w - z}\odif{z}
        = \dfrac{1}{2\pi
        i}\int_{\gamma}^{}\dfrac{f(w)(z - z_{0})^{n}}{(w - z_{0})^{n + 1}}\odif{w},\\
        &= \sum_{n = 0}^{\infty}\dfrac{(z - z_{0})^{n}}{2\pi
        i}\int_{\gamma}^{}\dfrac{f(w)}{(w
        - z_{0})^{n + 1}}\odif{w},\\
        \intertext{Por el la fórmula integral de Cacuchy generalizada,}
        &= \sum_{n = 0}^{\infty}\dfrac{(z - z_{0})^{n}}{n!}f^{(n)}(z_{0}),\\
        &=\sum_{n = 0}^{\infty} \dfrac{f^{(n)}(z_{0})}{n!}(z - z_{0})^{n}
    \end{align*}

    para toda \(z\) en \(\bar{B}_{r}(z_{0})\).
\end{proof}

\begin{example}
    Sea \(f(z) = \log(z + 1)\), entonces

    \begin{align*}
        f^{\prime}(z) &= \dfrac{1}{z + 1},\\
        f^{\prime\prime} &= - \dfrac{1}{(z + 1)^{2}},\\
        &\vdotswithin{ = }\\
        f^{(n)}(z) &= \dfrac{(n - 1)!(-1)^{n + 1}}{(z + 1)^{n}}.
    \end{align*}

    Por lo que

    \begin{align*}
        \log(z + 1) &= \sum_{n = 0}^{\infty}\dfrac{(n - 1)!(-1)^{n + 1}}{n!}z^{n},\\
        &= \sum_{n = 0}^{\infty}\dfrac{(-1)^{n + 1}}{n}z^{n}.\\
        \intertext{Y por el criterio de convergencia donde \(a_{n} = \tfrac{(-1)^{n +
        1}}{n}\),}
        &= \lim_{n \to \infty} \dfrac{\tfrac{1}{n}}{\tfrac{1}{n + 1}},\\
        &= \dfrac{n + 1}{n},\\
        &= 1.
    \end{align*}
\end{example}

\classdate{2026-05-20}

\begin{example}
    \begin{itemize}
        \item Sea \(f(z) = \mathrm{e}^{z}\) entera (analítica en
            \(\mathbb{C}\))

            Consideremos \(z_{0} = 0\). Sea \(n\in \mathbb{N},\; f^{(n)}(z) =
            \mathrm{e}^{z}\),
            tal que

            \begin{equation*}
                f^{(n)}(0) = \mathrm{e}^{0} = 1.
            \end{equation*}

            Y por Taylor,

            \begin{equation*}
                f(z) = \sum_{n = 0}^{\infty}\tfrac{1}{n!}z^{n},\; \forall z\in
                \bar{B}_{r}(0)\; \forall r > 0.
            \end{equation*}

            En particular, el radio de convergencia de

            \begin{equation*}
                \sum_{n = 0}^{\infty}\dfrac{1}{n!}z^{n}
            \end{equation*}

            es infinito.

        \item Sea \(\log\) la rama prinicipal del logaritmo (analítica en
                \(C\setminus (-\infty,
            0]\)). Definimos \(f(z) = \log(z + 1)\).

            Notemos que \(f\) es analítica en \(C\setminus (-\infty,
            -1] = U\). Sea \(0 < r < 1\). NOtemos que \(\bar{B}_{r}(0)
            \subseteq U\). Por Taylor,
            \(\forall z\in \bar{B}_{r}(0)\),

            \begin{equation*}
                f(z) = \sum_{n = 0}^{\infty}\dfrac{f^{(n)}(0)}{n!}z^{n}.
            \end{equation*}

            Entonces,

            \begin{align*}
                f^{\prime}(z) &= \dfrac{1}{z + 1},\\
                f^{\prime\prime}(z) &= - \dfrac{1}{(z + 1)^{2}},\\
                f^{(3)}(z) &= (-1)(2)\dfrac{1}{(z + 1)^{3}} =
                (-1)^{2}(2!)\dfrac{1}{(z + 1)^{3}},\\
                &\vdotswithin{ = }\\
                f^{(n)}(z) &= \dfrac{(-1)^{n - 1}(n - 1)!}{(z + 1)^{n}}.
            \end{align*}

            Y como esta serie está centrada en \(z_{0} = 0\),

            \begin{equation*}
                f^{(n)}(0) = (-1)^{n - 1}(n - 1)!.
            \end{equation*}

            Recordando que \(f^{(0)}(0) = \log(1)\), tal que

            \begin{equation*}
                \log(z + 1) = \sum_{n = 1}^{\infty}\dfrac{(n - 1)!(-1)^{n -
                1}}{n!}z^{n} = \sum_{n =
                1}^{\infty}\dfrac{(-1)^{n - 1}}{n}z^{n}\; \forall z\in
                \bar{B}_{r}(0)\; \forall r\in
                (0, 1).
            \end{equation*}

            Por lo que

            \begin{equation*}
                \log(z + 1) = \sum_{n = 1}^{\infty}\dfrac{(-1)^{n -
                1}}{n}z^{n}\; \forall z\in
                B_{1}(0).
            \end{equation*}

            \begin{remark}
                Si \(z = -1\), la serie es

                \begin{equation*}
                    -\sum_{n = 1}^{\infty}\dfrac{1}{n}
                \end{equation*}

                la cual diverge.
            \end{remark}

            Po lo tanto, por la observación \(R \ngtr 1\) pues diverge y
            tampoco puede ser menor a \(1\) por que contradice el hecho de
            que \(R\) esté dado en términos del supremo. Lo anterior  nos
            indica que el radio de convergencia de la serie es \(1\).
    \end{itemize}
\end{example}

\begin{lemma}
    \begin{equation*}
        (fg)^{n}(z) = \sum_{n = 0}^{\infty}
        \begin{pNiceMatrix}
            n\\
            k\\
        \end{pNiceMatrix}
        f^{k}(z)g^{n - k}(z).
    \end{equation*}
\end{lemma}

Veamos cómo se ven los coeficientes de la serie de Taylor del producto de funciones,

\begin{align*}
    \dfrac{(fg)^{(n)}(z_{0})}{n!} &= \dfrac{1}{n!}\sum_{k = 0}^{n}
    \begin{pNiceMatrix}
        n\\
        k\\
    \end{pNiceMatrix}
    f^{(k)}(z_{0})g^{(n - k)}(z_{0}),\\
    &= \sum_{k = 0}^{n}\dfrac{f^{(k)}(z_{0})}{k!}\dfrac{g^{(n - k)}(z_{0})}{(n -
    k)!} = c_{n}.
\end{align*}

Por lo tanto, la serie de Taylor del producto se ve como

\begin{equation*}
    (fg)(z) = \sum_{n = 0}^{\infty}(z - z_{0})^{n}.
\end{equation*}

\begin{example}
    Encuentra la serie de Taylor de \(\tfrac{z^{3}}{z - 1}\) en \(0\).

    Sea \(f(z) = z^{3}\) y \(g(z) = \tfrac{1}{z - 1}\). Entonces,

    \begin{equation*}
        f(z) = \sum_{n = 0}^{a_{n}z^{n}}\quad \wedge\quad g(z) = \sum_{n =
        0}^{\infty}b_{n}z^{n} = - \sum_{n = 0}^{\infty}z^{n}, \abs{z} < 1,
    \end{equation*}

    tal que

    \begin{equation*}
        a_{n} =
        \begin{dcases}
            1 & n = 3,\\
            0 & n \neq 3.
        \end{dcases}
        \quad
        \wedge
        \quad
        b_{n} = 1\; \forall n\in \mathbb{N}.
    \end{equation*}

    Por lo anterior,

    \begin{align*}
        \dfrac{z^{3}}{z - 1} &= \sum_{n = 0}^{\infty}\Bigl[\sum_{k = 0}^{n}a_{k}b_{n -
        k}\Bigr]z^{n},\\
        &= \sum_{n = 3}^{\infty}\Bigl[\sum_{k = 0}^{n}a_{k}(-1)\Bigr]z^{n},\\
        &= - \sum_{n = 3}^{\infty}z^{n}.
    \end{align*}
\end{example}

\classdate{2026-05-21}

\begin{lemma}[Dual de Abel]
    Sea \(\{b_{n}\} \subseteq \mathbb{C}\). Si

    \begin{equation*}
        \sum_{n = 1}^{\infty}\dfrac{b_{n}}{(z - z_{0})^{n}}
    \end{equation*}

    converge en \(z_{1}\in \mathbb{C}\), entonces la serie converge absolutamente en
    \(\set{z\in \mathbb{C} \given \abs{z - z_{0}} > r_{1}}\) donde \(r_{1} =
        \abs{z_{1} -
    z_{0}}\). Además, converge uniformemente en \(\set{z\in \mathbb{C} \given \abs{z -
    z_{0}} \geq r}\; \forall r > r_{1}\).
\end{lemma}

\begin{proof}
    Como la serie

    \begin{equation*}
        \sum_{n = 1}^{\infty}\dfrac{b_{n}}{(z - z_{0})^{n}}
    \end{equation*}

    converge en \(z_{1}\), la sucesión \(\{\tfrac{\abs{b_{n}}}{r_{1}^{n}}\}\)
    es acotada, es decir, \(\exists M > 0\) tal que \(\forall n\in
    \mathbb{N},\; \tfrac{\abs{b_{n}}}{r_{1}^{n}} < M\). Por lo que basta
    mostrar que si \(r > r_{1}\), la serie

    \begin{equation*}
        \sum_{n = 1}^{\infty}\dfrac{b_{n}}{(z - z_{0})^{n}}
    \end{equation*}

    converge uniforme y absolutamente en \(\set{z\in \mathbb{C} \given \abs{z
    - z_{0}} \geq r}\),

    \begin{equation*}
        \abs`\bigg{\dfrac{b_{n}}{(z - z_{0})^{n}}} \leq \dfrac{\abs{b_{n}}}{r^{n}} =
        \dfrac{\abs{b_{n}}}{r_{1}^{n}}\dfrac{r_{1}^{n}}{r^{n}} \leq M
        \bigl(\underbrace{\dfrac{r_{1}}{r}}_{\in (0, 1)}\bigr)^{n}.
    \end{equation*}

    Por lo tanto,

    \begin{equation*}
        \sum_{n = 0}^{\infty}\bigl(\dfrac{r_{1}}{r}\bigr)^{n}
    \end{equation*}

    converge.

    Y por la prueba M de Weierstrass converge uniforme y absolutamente.
\end{proof}

\begin{lemma}
    Sea \(f\) una función analítica en el anillo

    \begin{equation*}
        A_{r_{1}}^{r_{2}} = \set{z\in \mathbb{C} \given r_{1} < \abs{z - z_{0}} < r_{2}}
    \end{equation*}

    con \(r_{1} < r_{2}\) y \(\gamma_{j}(\theta) = z_{0}t +
    t_{j}\mathrm{e}^{i\theta}\), \(j = {1, 2},\; r_{1} < t_{1} < t_{2} <
    r_{2},\; \theta\in [0, 2\pi]\). Entonces \(\forall z\in A_{t_{1}}^{t_{2}}\),

    \begin{equation*}
        f(z) = \dfrac{1}{2\pi i}\int_{\gamma_{2}}^{}\dfrac{f(w)}{w - z}\odif{w} -
        \dfrac{1}{2\pi i}\int_{\gamma_{1}}^{}\dfrac{f(w)}{w - z}\odif{w}.
    \end{equation*}
\end{lemma}

\begin{proof}
    Sea \(z\in A_{t_{1}}^{t_{2}}\).
    Definimos la función auxiliar

    \begin{equation*}
        g(w) =
        \begin{dcases}
            \dfrac{f(w) - f(z)}{w - z}, & w \neq z,\\
            f^{\prime}(z), & w = z
        \end{dcases}
    \end{equation*}

    analítica en \(A\).

    Y como \(\gamma_{1}\) es homotópica a \(\gamma_{2}\) en \(A\),

    \begin{equation*}
        \int_{\gamma_{2}}^{}g(w)\odif{w} - \int_{\gamma_{1}}^{}g(w)\odif{w} = 0.
    \end{equation*}

    Por lo que,

    \begin{equation*}
        \int_{\gamma_{2}}^{}\dfrac{f(w)}{w - z}\odif{w} -
        f(z)\int_{\gamma_{2}}^{}\dfrac{1}{w
        - 1}\odif{w} - \int_{\gamma_{1}}^{}\dfrac{f(w)}{w - z}\odif{w} + f(z)
        \int_{\gamma_{1}}^{}\dfrac{1}{w - z}\odif{w} = 0,
    \end{equation*}

    donde la última integral es \(0\) por el teorema de Cauchy y por el
    número de vueltas
    \(n(\gamma_{2}, z)\),

    \begin{equation*}
        -f(z) = \int_{\gamma_{2}}^{}\dfrac{1}{w - z}\odif{w} = -f(z)2\pi i.
    \end{equation*}

    Por lo tanto,

    \begin{equation*}
        f(z) = \dfrac{1}{2\pi i}\int_{\gamma_{2}}^{}\dfrac{f(w)}{w - z}\odif{w} -
        \dfrac{1}{2\pi i}\int_{\gamma_{1}}^{}\dfrac{f(w)}{w - z}\odif{w}.
    \end{equation*}
\end{proof}

\begin{theorem}[de Laurent]
    Sea \(f\) analítica en \(A_{r_{1}}^{r_{2}}\). Entonces \(\forall z\in
    A_{r_{1}}^{r_{2}}\),

    \begin{equation*}
        f(z) = \sum_{n = 0}^{\infty}a_{n}(z - z_{0})^{n} + \sum_{n =
        1}^{\infty}\dfrac{b_{n}}{(z - z_{0})^{n}}
    \end{equation*}

    converge absolutamente.

    A

    \begin{equation*}
        \sum_{n = 0}^{\infty}a_{n}(z - z_{0})^{n} + \sum_{n =
        1}^{\infty}\dfrac{b_{n}}{(z - z_{0})^{n}}
    \end{equation*}

    se le conce como serie de Laurent y en ocasiones por ``abuso de
    notación'' se escribe
    como

    \begin{equation*}
        \sum_{n = -\infty}^{\infty}c_{n}(z - z_{0})^{n},
    \end{equation*}

    donde

    \begin{equation*}
        c_{n} =
        \begin{dcases}
            a_{n} & n \geq 0,\\
            b_{n} & n < 0.
        \end{dcases}
    \end{equation*}

    Además la convergencia uniforme en subanillos

    \begin{equation*}
        A_{s_{1}}^{s_{2}} = \set{z\in \mathbb{C} \given s_{1} \leq \abs{z -
        z_{0}} \leq s_{2}}
    \end{equation*}

    con \(r_{1} < s_{1} < s_{2} < r_{2}\).
\end{theorem}

Antes de continar con la prueba del teorema, veamos un ejemplo.

\begin{example}
    Encuentre la serie de Laurent \(f(z) = \tfrac{1}{z^{2} - z^{3}}\) en \(0\).

    Notemos que \(f(z) = \tfrac{1}{1 - z}\tfrac{1}{z^{2}}\) y para \(\abs{z} < 1\),
    \(\tfrac{1}{1 - z} = \sum_{n = 0}^{\infty}z^{n}\). Por lo tanto,

    \begin{equation*}
        \dfrac{1}{1 - z}\dfrac{1}{z^{2}} = \dfrac{1}{z^{2}} + \dfrac{1}{z} +
        1 + z + z^{2} +
        \cdots
    \end{equation*}
\end{example}
\end{document}
